```latex
\documentclass[12pt,a4paper]{article}

% ==================================================
% PACKAGE
% ==================================================
\usepackage[utf8]{inputenc}
\usepackage[T1]{fontenc}
\usepackage[indonesian]{babel}
\usepackage{geometry}
\usepackage{setspace}
\usepackage{parskip}
\usepackage{titlesec}
\usepackage{xcolor}
\usepackage{listings}
\usepackage{hyperref}
\usepackage{enumitem}
\usepackage{fancyhdr}
\usepackage{tabularx}
\usepackage{booktabs}
\usepackage{tcolorbox}

% ==================================================
% PAGE CONFIGURATION
% ==================================================
\geometry{
    left=3cm,
    right=2.5cm,
    top=2.5cm,
    bottom=2.5cm
}

\setstretch{1.2}
\setlength{\parindent}{0pt}
\setlength{\parskip}{6pt}

% ==================================================
% COLOR
% ==================================================
\definecolor{dawetblue}{RGB}{25,55,95}
\definecolor{codebg}{RGB}{245,247,250}
\definecolor{codeborder}{RGB}{210,218,230}
\definecolor{softblue}{RGB}{235,242,250}
\definecolor{darkgray}{RGB}{55,65,75}
\definecolor{warningyellow}{RGB}{255,247,220}
\definecolor{successgreen}{RGB}{225,243,232}

% ==================================================
% HYPERLINK
% ==================================================
\hypersetup{
    colorlinks=true,
    linkcolor=dawetblue,
    urlcolor=dawetblue,
    citecolor=dawetblue
}

% ==================================================
% SECTION STYLE
% ==================================================
\titleformat{\section}
{\Large\bfseries\color{dawetblue}}
{\thesection.}{0.6em}{}

\titleformat{\subsection}
{\large\bfseries\color{darkgray}}
{\thesubsection}{0.6em}{}

% ==================================================
% CODE LISTING
% ==================================================
\lstdefinestyle{dawetcode}{
    backgroundcolor=\color{codebg},
    basicstyle=\ttfamily\small,
    frame=single,
    rulecolor=\color{codeborder},
    breaklines=true,
    columns=fullflexible,
    keepspaces=true,
    showstringspaces=false,
    tabsize=2,
    numbers=none,
    xleftmargin=8pt,
    xrightmargin=8pt,
    aboveskip=8pt,
    belowskip=8pt
}

\lstset{style=dawetcode}

% ==================================================
% HEADER & FOOTER
% ==================================================
\pagestyle{fancy}
\fancyhf{}
\lhead{\small PROJECT DAWET}
\rhead{\small Minggu 8}
\cfoot{\thepage}
\renewcommand{\headrulewidth}{0.4pt}

% ==================================================
% DOCUMENT
% ==================================================
\begin{document}

% ==================================================
% COVER
% ==================================================
\begin{titlepage}
    \centering

    \vspace*{2cm}

    {\Large\bfseries PROJECT DAWET\par}

    \vspace{1.5cm}

    {\Huge\bfseries JavaScript Fundamentals\par}

    \vspace{0.5cm}

    {\Large Minggu 8 -- Assignment 2\par}

    \vspace{1.5cm}

    \begin{tcolorbox}[
        colback=softblue,
        colframe=dawetblue,
        width=0.85\textwidth,
        arc=2mm,
        boxrule=0.8pt
    ]
        \centering
        \textbf{Learning Focus}

        Memahami dasar-dasar pemrograman JavaScript,
        mulai dari variabel, tipe data, operator,
        percabangan, perulangan, fungsi, array,
        hingga object.
    \end{tcolorbox}

    \vfill

    {\large Asynchronous Learning Program\par}

    \vspace{0.3cm}

    {\large Learn. Build. Understand.\par}

    \vspace{1cm}
\end{titlepage}

\tableofcontents
\newpage

% ==================================================
\section{Pendahuluan}
% ==================================================

Pada materi sebelumnya, kita telah mempelajari HTML
untuk membangun struktur halaman web dan CSS untuk
mengatur tampilan serta tata letaknya.

Namun, sebuah website tidak selalu cukup hanya
memiliki struktur dan tampilan. Dalam banyak kasus,
website juga membutuhkan kemampuan untuk merespons
tindakan pengguna, memproses data, dan menjalankan
logika tertentu.

Contohnya:

\begin{itemize}
    \item Menghitung total harga produk.
    \item Menampilkan pesan ketika tombol ditekan.
    \item Memeriksa apakah formulir sudah diisi.
    \item Menampilkan daftar produk berdasarkan data.
    \item Menghitung jumlah barang yang dipilih pengguna.
\end{itemize}

Untuk membuat perilaku tersebut, kita membutuhkan
JavaScript.

JavaScript merupakan bahasa pemrograman yang banyak
digunakan dalam pengembangan web. Pada sisi browser,
JavaScript dapat digunakan untuk membuat halaman web
menjadi interaktif.

Dalam rangkaian materi Assignment 2, JavaScript juga
akan menjadi dasar untuk mempelajari manipulasi DOM,
pengolahan formulir, serta komunikasi dengan backend.

\subsection{Tujuan Pembelajaran}

Setelah menyelesaikan materi ini, peserta diharapkan
mampu:

\begin{enumerate}
    \item Menjelaskan fungsi JavaScript dalam pengembangan web.
    \item Menghubungkan file JavaScript dengan HTML.
    \item Menggunakan variabel dan tipe data.
    \item Menggunakan operator dalam ekspresi.
    \item Membuat percabangan berdasarkan kondisi.
    \item Menggunakan perulangan untuk memproses data.
    \item Membuat dan memanggil fungsi.
    \item Menggunakan array dan object sederhana.
    \item Menyelesaikan permasalahan sederhana menggunakan
    logika JavaScript.
\end{enumerate}

% ==================================================
\section{Mengenal JavaScript}
% ==================================================

\subsection{Apa Itu JavaScript?}

JavaScript adalah bahasa pemrograman yang digunakan
untuk membuat aplikasi interaktif.

Dalam pengembangan web, JavaScript dapat berjalan
di browser pengguna dan berinteraksi dengan halaman
HTML melalui Document Object Model (DOM).

Contoh penerapannya:

\begin{itemize}
    \item Tombol yang menjalankan suatu aksi.
    \item Formulir yang memeriksa input pengguna.
    \item Menu navigasi interaktif.
    \item Perhitungan harga secara otomatis.
    \item Pengambilan dan penampilan data dari server.
\end{itemize}

JavaScript juga dapat digunakan di luar browser,
misalnya melalui Node.js untuk membangun aplikasi
backend.

Namun, pada materi ini kita akan memulai dari
penggunaan JavaScript di browser.

\subsection{Hubungan HTML, CSS, dan JavaScript}

Ketiganya memiliki peran yang berbeda, tetapi saling
melengkapi dalam membangun sebuah website.

\begin{center}
\begin{tabularx}{\textwidth}{lX}
\toprule
\textbf{Teknologi} & \textbf{Peran Utama} \\
\midrule
HTML & Membentuk struktur dan konten halaman. \\
CSS & Mengatur tampilan, warna, ukuran, dan tata letak. \\
JavaScript & Menambahkan logika dan perilaku interaktif. \\
\bottomrule
\end{tabularx}
\end{center}

Sebagai contoh, sebuah tombol pada website dapat
dibangun menggunakan HTML, diberi warna menggunakan
CSS, kemudian diberi fungsi tertentu menggunakan
JavaScript.

\subsection{Cara Menambahkan JavaScript ke HTML}

Terdapat beberapa cara untuk menggunakan JavaScript
dalam dokumen HTML.

\subsubsection{Inline JavaScript}

JavaScript dapat ditulis langsung pada atribut HTML
tertentu.

Contoh:

\begin{lstlisting}
<button onclick="alert('Halo DAWET!')">
    Klik Saya
</button>
\end{lstlisting}

Cara ini mudah digunakan untuk demonstrasi sederhana,
tetapi tidak disarankan untuk pengembangan proyek
yang lebih terstruktur karena logika JavaScript
tercampur dengan HTML.

\subsubsection{Internal JavaScript}

JavaScript dapat ditulis menggunakan elemen
\texttt{<script>} di dalam file HTML.

Contoh:

\begin{lstlisting}
<!DOCTYPE html>
<html lang="id">
<head>
    <meta charset="UTF-8">
    <title>Belajar JavaScript</title>
</head>
<body>

    <h1>Belajar JavaScript</h1>

    <script>
        console.log("Halo DAWET!");
    </script>

</body>
</html>
\end{lstlisting}

\subsubsection{External JavaScript}

Cara yang direkomendasikan untuk proyek yang
memiliki banyak kode JavaScript adalah memisahkan
kode ke dalam file tersendiri.

Contoh struktur proyek:

\begin{lstlisting}
belajar-javascript/
|
|-- index.html
|-- script.js
\end{lstlisting}

Isi file \texttt{index.html}:

\begin{lstlisting}
<!DOCTYPE html>
<html lang="id">
<head>
    <meta charset="UTF-8">
    <meta name="viewport"
          content="width=device-width, initial-scale=1.0">
    <title>Belajar JavaScript</title>
</head>
<body>

    <h1>Belajar JavaScript</h1>

    <script src="script.js"></script>
</body>
</html>
\end{lstlisting}

Isi file \texttt{script.js}:

\begin{lstlisting}
console.log("Halo DAWET!");
\end{lstlisting}

Atribut \texttt{src} digunakan untuk menentukan
lokasi file JavaScript yang ingin dihubungkan.

Pada contoh tersebut, file JavaScript berada
dalam direktori yang sama dengan file HTML.

Untuk proyek ini, kita akan menggunakan external
JavaScript agar kode lebih terorganisasi.

\subsection{Mengenal Console}

Console merupakan fitur pada browser yang dapat
digunakan untuk menampilkan informasi dan membantu
proses debugging.

Salah satu perintah yang sering digunakan adalah:

\begin{lstlisting}
console.log("Halo, dunia!");
\end{lstlisting}

Untuk membuka console di browser:

\begin{enumerate}
    \item Buka halaman HTML melalui browser.
    \item Tekan tombol F12 atau klik kanan lalu pilih
    \textit{Inspect}.
    \item Pilih tab \textit{Console}.
    \item Perhatikan keluaran dari perintah
    \texttt{console.log()}.
\end{enumerate}

Console akan menjadi alat penting selama proses
belajar JavaScript.

% ==================================================
\section{Variabel dan Tipe Data}
% ==================================================

\subsection{Apa Itu Variabel?}

Variabel merupakan tempat untuk menyimpan nilai
yang dapat digunakan kembali dalam program.

Contohnya, sebuah website UMKM mungkin perlu
menyimpan informasi nama produk, harga, dan jumlah
stok.

Dalam JavaScript, variabel dapat dibuat menggunakan
kata kunci \texttt{let}, \texttt{const}, atau
\texttt{var}.

\subsection{Deklarasi Variabel}

\textbf{1. let}

Digunakan untuk membuat variabel yang nilainya
dapat diubah.

\begin{lstlisting}
let namaProduk = "Kopi Lokal";

namaProduk = "Kopi Susu";

console.log(namaProduk);
\end{lstlisting}

Hasil:

\begin{lstlisting}
Kopi Susu
\end{lstlisting}

\textbf{2. const}

Digunakan untuk membuat variabel yang tidak dapat
diisi ulang melalui assignment setelah deklarasi.

\begin{lstlisting}
const namaToko = "DAWET Store";

console.log(namaToko);
\end{lstlisting}

Contoh berikut akan menghasilkan error:

\begin{lstlisting}
const harga = 15000;
harga = 20000;
\end{lstlisting}

Perlu diketahui bahwa \texttt{const} tidak selalu
berarti seluruh isi objek atau array menjadi
tidak dapat diubah. Yang tidak dapat diubah adalah
pengikatan variabelnya ke nilai lain.

\textbf{3. var}

\texttt{var} merupakan cara deklarasi variabel
yang digunakan pada JavaScript versi lama.

\begin{lstlisting}
var nama = "Budi";
console.log(nama);
\end{lstlisting}

Pada pembelajaran ini, utamakan penggunaan
\texttt{let} dan \texttt{const} karena keduanya
memiliki aturan cakupan variabel yang lebih mudah
dikelola untuk kode modern.

\subsection{Aturan Penamaan Variabel}

Beberapa aturan dasar penamaan variabel:

\begin{itemize}
    \item Nama variabel dapat menggunakan huruf,
    angka, garis bawah, dan tanda dolar.
    \item Nama variabel tidak boleh diawali angka.
    \item Nama variabel tidak boleh menggunakan
    kata kunci JavaScript.
    \item JavaScript membedakan huruf besar dan kecil.
\end{itemize}

Contoh:

\begin{lstlisting}
let namaProduk = "Nasi Goreng";
let hargaProduk = 15000;
let jumlahStok = 20;
\end{lstlisting}

Gunakan nama yang menjelaskan isi variabel agar
kode lebih mudah dibaca.

\subsection{Tipe Data}

Tipe data menjelaskan jenis nilai yang disimpan
dalam sebuah variabel.

Beberapa tipe data dasar JavaScript adalah:

\textbf{1. String}

Digunakan untuk menyimpan teks.

\begin{lstlisting}
let nama = "Kopi Susu";
let kategori = 'Minuman';
\end{lstlisting}

\textbf{2. Number}

Digunakan untuk menyimpan angka, baik bilangan
bulat maupun desimal.

\begin{lstlisting}
let harga = 15000;
let rating = 4.8;
\end{lstlisting}

\textbf{3. Boolean}

Digunakan untuk menyimpan nilai logika
\texttt{true} atau \texttt{false}.

\begin{lstlisting}
let tersedia = true;
let sudahDibayar = false;
\end{lstlisting}

\textbf{4. Undefined}

Menunjukkan variabel yang telah dideklarasikan,
tetapi belum diberi nilai.

\begin{lstlisting}
let alamat;

console.log(alamat);
\end{lstlisting}

\textbf{5. Null}

Digunakan untuk menyatakan bahwa suatu nilai
sengaja dikosongkan.

\begin{lstlisting}
let nomorTelepon = null;
\end{lstlisting}

\textbf{6. Object}

Digunakan untuk mengelompokkan data dalam bentuk
pasangan properti dan nilai.

\begin{lstlisting}
const produk = {
    nama: "Kopi Susu",
    harga: 15000,
    tersedia: true
};
\end{lstlisting}

\textbf{7. Array}

Digunakan untuk menyimpan kumpulan nilai dalam
satu variabel.

\begin{lstlisting}
const daftarProduk = [
    "Kopi Susu",
    "Teh Manis",
    "Roti Bakar"
];
\end{lstlisting}

\subsection{Memeriksa Tipe Data}

Gunakan operator \texttt{typeof} untuk memeriksa
tipe data suatu nilai.

\begin{lstlisting}
let nama = "DAWET";
let harga = 15000;
let aktif = true;

console.log(typeof nama);
console.log(typeof harga);
console.log(typeof aktif);
\end{lstlisting}

Hasil:

\begin{lstlisting}
string
number
boolean
\end{lstlisting}

Perlu diperhatikan bahwa \texttt{typeof null}
menghasilkan \texttt{"object"} karena perilaku
historis JavaScript.

% ==================================================
\section{Operator dalam JavaScript}
% ==================================================

Operator digunakan untuk melakukan operasi terhadap
nilai atau variabel.

\subsection{Operator Aritmatika}

Operator aritmatika digunakan untuk melakukan
perhitungan matematika.

\begin{center}
\begin{tabularx}{\textwidth}{cX}
\toprule
\textbf{Operator} & \textbf{Fungsi} \\
\midrule
+ & Penjumlahan \\
- & Pengurangan \\
* & Perkalian \\
/ & Pembagian \\
\% & Sisa hasil bagi atau modulo \\
** & Perpangkatan \\
\bottomrule
\end{tabularx}
\end{center}

Contoh:

\begin{lstlisting}
let harga = 15000;
let jumlah = 3;

let total = harga * jumlah;

console.log(total);
\end{lstlisting}

Hasil:

\begin{lstlisting}
45000
\end{lstlisting}

\subsection{Operator Perbandingan}

Operator perbandingan digunakan untuk membandingkan
dua nilai dan menghasilkan nilai Boolean.

\begin{center}
\begin{tabularx}{\textwidth}{cX}
\toprule
\textbf{Operator} & \textbf{Fungsi} \\
\midrule
== & Membandingkan nilai dengan konversi tipe tertentu. \\
=== & Membandingkan nilai dan tipe data. \\
!= & Membandingkan ketidaksamaan nilai. \\
!== & Membandingkan ketidaksamaan nilai atau tipe. \\
> & Lebih besar dari. \\
< & Lebih kecil dari. \\
>= & Lebih besar atau sama dengan. \\
<= & Lebih kecil atau sama dengan. \\
\bottomrule
\end{tabularx}
\end{center}

Contoh:

\begin{lstlisting}
let harga = 15000;

console.log(harga > 10000);
console.log(harga === 15000);
\end{lstlisting}

Hasil:

\begin{lstlisting}
true
true
\end{lstlisting}

Dalam praktik sehari-hari, gunakan
\texttt{===} dan \texttt{!==} sebagai pilihan
utama agar perbandingan lebih ketat dan mudah
diprediksi.

\subsection{Operator Logika}

Operator logika digunakan untuk menggabungkan
atau membalikkan kondisi.

\begin{center}
\begin{tabularx}{\textwidth}{cX}
\toprule
\textbf{Operator} & \textbf{Fungsi} \\
\midrule
\&\& & AND: semua kondisi harus bernilai benar. \\
|| & OR: minimal satu kondisi harus benar. \\
! & NOT: membalikkan nilai logika. \\
\bottomrule
\end{tabularx}
\end{center}

Contoh AND:

\begin{lstlisting}
let stok = 10;
let aktif = true;

console.log(stok > 0 && aktif);
\end{lstlisting}

Hasil:

\begin{lstlisting}
true
\end{lstlisting}

Contoh OR:

\begin{lstlisting}
let metodePembayaran = "transfer";

console.log(
    metodePembayaran === "transfer" ||
    metodePembayaran === "tunai"
);
\end{lstlisting}

Hasil:

\begin{lstlisting}
true
\end{lstlisting}

Contoh NOT:

\begin{lstlisting}
let tersedia = true;

console.log(!tersedia);
\end{lstlisting}

Hasil:

\begin{lstlisting}
false
\end{lstlisting}

% ==================================================
\section{Percabangan}
% ==================================================

\subsection{Apa Itu Percabangan?}

Percabangan adalah mekanisme yang memungkinkan
program mengambil keputusan berdasarkan kondisi
tertentu.

Sebagai contoh, sebuah sistem UMKM dapat menentukan
apakah suatu produk masih tersedia berdasarkan
jumlah stoknya.

\subsection{if}

Struktur \texttt{if} menjalankan kode ketika
kondisi bernilai benar.

\begin{lstlisting}
let stok = 10;

if (stok > 0) {
    console.log("Produk tersedia");
}
\end{lstlisting}

Karena stok lebih besar dari nol, pesan
\texttt{Produk tersedia} akan ditampilkan.

\subsection{if dan else}

Tambahkan \texttt{else} jika program perlu
menjalankan alternatif ketika kondisi salah.

\begin{lstlisting}
let stok = 0;

if (stok > 0) {
    console.log("Produk tersedia");
} else {
    console.log("Produk habis");
}
\end{lstlisting}

\subsection{if, else if, dan else}

Struktur ini digunakan ketika terdapat beberapa
kemungkinan kondisi.

\begin{lstlisting}
let totalBelanja = 75000;

if (totalBelanja >= 100000) {
    console.log("Diskon 20%");
} else if (totalBelanja >= 50000) {
    console.log("Diskon 10%");
} else {
    console.log("Tidak ada diskon");
}
\end{lstlisting}

Program akan memeriksa kondisi dari atas ke bawah.
Ketika suatu kondisi terpenuhi, blok yang sesuai
dijalankan dan rangkaian percabangan tersebut
selesai.

\subsection{switch}

\texttt{switch} dapat digunakan ketika kita ingin
memeriksa satu nilai terhadap beberapa kemungkinan
nilai yang berbeda.

\begin{lstlisting}
let kategori = "minuman";

switch (kategori) {
    case "makanan":
        console.log("Kategori makanan");
        break;

    case "minuman":
        console.log("Kategori minuman");
        break;

    case "camilan":
        console.log("Kategori camilan");
        break;

    default:
        console.log("Kategori tidak dikenal");
}
\end{lstlisting}

Perintah \texttt{break} digunakan untuk menghentikan
eksekusi pada case yang telah cocok agar program
tidak melanjutkan ke case berikutnya.

\texttt{default} dijalankan jika tidak ada case
yang cocok.

% ==================================================
\section{Perulangan}
% ==================================================

\subsection{Apa Itu Perulangan?}

Perulangan atau loop digunakan untuk menjalankan
blok kode berulang kali.

Misalnya, kita ingin menampilkan daftar sepuluh
produk. Daripada menulis perintah yang sama
sepuluh kali, kita dapat menggunakan perulangan.

\subsection{for}

Perulangan \texttt{for} biasanya digunakan ketika
jumlah pengulangan diketahui atau dapat ditentukan.

Struktur dasarnya:

\begin{lstlisting}
for (inisialisasi; kondisi; perubahan) {
    // kode yang dijalankan
}
\end{lstlisting}

Contoh:

\begin{lstlisting}
for (let i = 1; i <= 5; i++) {
    console.log("Produk ke-" + i);
}
\end{lstlisting}

Penjelasan:

\begin{itemize}
    \item \texttt{let i = 1}: nilai awal penghitung.
    \item \texttt{i <= 5}: perulangan berjalan selama
    kondisi benar.
    \item \texttt{i++}: menaikkan nilai penghitung
    sebanyak satu pada setiap putaran.
\end{itemize}

\subsection{while}

Perulangan \texttt{while} menjalankan kode selama
kondisi bernilai benar.

\begin{lstlisting}
let stok = 3;

while (stok > 0) {
    console.log("Stok tersisa: " + stok);
    stok--;
}
\end{lstlisting}

Perulangan berhenti ketika nilai \texttt{stok}
tidak lagi lebih besar dari nol.

\begin{tcolorbox}[
    colback=warningyellow,
    colframe=orange!70!black,
    title=\textbf{Perhatian: Infinite Loop}
]
Pastikan kondisi perulangan dapat berubah sehingga
pada akhirnya menjadi salah.

Jika kondisi selalu benar dan tidak ada mekanisme
penghentian, program dapat mengalami infinite loop,
yaitu perulangan yang tidak pernah selesai.
\end{tcolorbox}

\subsection{break dan continue}

\texttt{break} digunakan untuk menghentikan
perulangan secara langsung.

\begin{lstlisting}
for (let i = 1; i <= 5; i++) {
    if (i === 3) {
        break;
    }

    console.log(i);
}
\end{lstlisting}

Hasil:

\begin{lstlisting}
1
2
\end{lstlisting}

Sementara itu, \texttt{continue} digunakan untuk
melewati satu putaran dan melanjutkan ke putaran
berikutnya.

\begin{lstlisting}
for (let i = 1; i <= 5; i++) {
    if (i === 3) {
        continue;
    }

    console.log(i);
}
\end{lstlisting}

Hasil:

\begin{lstlisting}
1
2
4
5
\end{lstlisting}

% ==================================================
\section{Fungsi}
% ==================================================

\subsection{Apa Itu Fungsi?}

Fungsi adalah blok kode yang dirancang untuk
menjalankan tugas tertentu dan dapat dipanggil
kembali ketika dibutuhkan.

Fungsi membantu kita menghindari penulisan kode
yang sama secara berulang.

Sebagai contoh, jika kita sering menghitung total
harga produk, kita dapat membuat satu fungsi
untuk menangani perhitungan tersebut.

\subsection{Membuat Function Declaration}

Contoh:

\begin{lstlisting}
function sapa() {
    console.log("Halo, selamat datang!");
}
\end{lstlisting}

Kode di atas hanya mendefinisikan fungsi.
Agar dijalankan, kita perlu memanggilnya:

\begin{lstlisting}
sapa();
\end{lstlisting}

\subsection{Parameter dan Argument}

Fungsi dapat menerima data melalui parameter.

\begin{lstlisting}
function sapa(nama) {
    console.log("Halo, " + nama);
}

sapa("Budi");
\end{lstlisting}

Pada contoh tersebut:

\begin{itemize}
    \item \texttt{nama} adalah parameter.
    \item \texttt{"Budi"} adalah argument yang dikirim
    ketika fungsi dipanggil.
\end{itemize}

\subsection{return}

Fungsi dapat menghasilkan nilai menggunakan
perintah \texttt{return}.

\begin{lstlisting}
function hitungTotal(harga, jumlah) {
    return harga * jumlah;
}

let total = hitungTotal(15000, 3);

console.log(total);
\end{lstlisting}

Hasil:

\begin{lstlisting}
45000
\end{lstlisting}

Perbedaan penting:

\begin{itemize}
    \item \texttt{console.log()} menampilkan informasi
    pada console.
    \item \texttt{return} mengembalikan nilai dari fungsi
    kepada bagian program yang memanggilnya.
\end{itemize}

\subsection{Function Expression}

Fungsi juga dapat disimpan ke dalam variabel.

\begin{lstlisting}
const hitungLuas = function(panjang, lebar) {
    return panjang * lebar;
};

console.log(hitungLuas(5, 4));
\end{lstlisting}

\subsection{Arrow Function}

JavaScript modern menyediakan bentuk penulisan
fungsi yang disebut arrow function.

\begin{lstlisting}
const hitungTotal = (harga, jumlah) => {
    return harga * jumlah;
};

console.log(hitungTotal(15000, 3));
\end{lstlisting}

Untuk fungsi dengan satu ekspresi sederhana,
penulisannya dapat dibuat lebih singkat:

\begin{lstlisting}
const tambah = (a, b) => a + b;

console.log(tambah(5, 3));
\end{lstlisting}

Pada tahap awal, peserta cukup memahami cara
mendefinisikan fungsi, memberikan parameter,
memanggil fungsi, dan mengembalikan nilai.

% ==================================================
\section{Array}
% ==================================================

\subsection{Apa Itu Array?}

Array adalah struktur data yang digunakan untuk
menyimpan kumpulan nilai dalam satu variabel.

Contoh:

\begin{lstlisting}
const produk = [
    "Kopi Susu",
    "Teh Manis",
    "Roti Bakar"
];
\end{lstlisting}

Setiap nilai di dalam array memiliki posisi
yang disebut index.

Index array JavaScript dimulai dari angka nol.

\begin{center}
\begin{tabular}{c|c|c|c}
\toprule
Index & 0 & 1 & 2 \\
\midrule
Nilai & Kopi Susu & Teh Manis & Roti Bakar \\
\bottomrule
\end{tabular}
\end{center}

\subsection{Mengakses Elemen Array}

Gunakan index untuk mengambil elemen tertentu.

\begin{lstlisting}
const produk = [
    "Kopi Susu",
    "Teh Manis",
    "Roti Bakar"
];

console.log(produk[0]);
console.log(produk[1]);
\end{lstlisting}

Hasil:

\begin{lstlisting}
Kopi Susu
Teh Manis
\end{lstlisting}

\subsection{Mengubah Elemen Array}

Elemen array dapat diubah melalui index.

\begin{lstlisting}
const produk = [
    "Kopi Susu",
    "Teh Manis"
];

produk[1] = "Es Teh";

console.log(produk);
\end{lstlisting}

Meskipun menggunakan \texttt{const}, elemen
array tetap dapat diubah. Yang tidak dapat
diubah adalah pengikatan variabel \texttt{produk}
ke array lain.

\subsection{Properti length}

Properti \texttt{length} digunakan untuk mengetahui
jumlah elemen dalam array.

\begin{lstlisting}
const produk = [
    "Kopi Susu",
    "Teh Manis",
    "Roti Bakar"
];

console.log(produk.length);
\end{lstlisting}

Hasil:

\begin{lstlisting}
3
\end{lstlisting}

\subsection{Method Dasar Array}

Beberapa method yang umum digunakan:

\begin{center}
\begin{tabularx}{\textwidth}{lX}
\toprule
\textbf{Method} & \textbf{Fungsi} \\
\midrule
push() & Menambahkan elemen di akhir array. \\
pop() & Menghapus elemen terakhir. \\
unshift() & Menambahkan elemen di awal array. \\
shift() & Menghapus elemen pertama. \\
includes() & Memeriksa apakah nilai terdapat dalam array. \\
indexOf() & Mencari posisi pertama suatu nilai. \\
\bottomrule
\end{tabularx}
\end{center}

Contoh:

\begin{lstlisting}
const produk = ["Kopi", "Teh"];

produk.push("Roti");

console.log(produk);
\end{lstlisting}

Hasil:

\begin{lstlisting}
["Kopi", "Teh", "Roti"]
\end{lstlisting}

\subsection{Menggunakan Perulangan pada Array}

Array dapat diproses menggunakan perulangan.

\begin{lstlisting}
const produk = [
    "Kopi Susu",
    "Teh Manis",
    "Roti Bakar"
];

for (let i = 0; i < produk.length; i++) {
    console.log(produk[i]);
}
\end{lstlisting}

Kita juga dapat menggunakan \texttt{for...of}
untuk membaca setiap nilai array.

\begin{lstlisting}
for (const namaProduk of produk) {
    console.log(namaProduk);
}
\end{lstlisting}

% ==================================================
\section{Object}
% ==================================================

\subsection{Apa Itu Object?}

Object merupakan struktur data yang digunakan
untuk menyimpan kumpulan informasi dalam bentuk
pasangan key dan value.

Sebagai contoh, sebuah produk UMKM memiliki
nama, harga, kategori, dan status ketersediaan.

Informasi tersebut dapat dikelompokkan dalam
satu object.

\begin{lstlisting}
const produk = {
    nama: "Kopi Susu",
    harga: 15000,
    kategori: "Minuman",
    tersedia: true
};
\end{lstlisting}

Pada contoh tersebut:

\begin{itemize}
    \item \texttt{nama} adalah key dengan nilai
    \texttt{"Kopi Susu"}.
    \item \texttt{harga} adalah key dengan nilai
    \texttt{15000}.
    \item \texttt{kategori} adalah key dengan nilai
    \texttt{"Minuman"}.
    \item \texttt{tersedia} adalah key dengan nilai
    \texttt{true}.
\end{itemize}

\subsection{Mengakses Properti Object}

Terdapat dua cara umum untuk mengakses properti
object.

Menggunakan dot notation:

\begin{lstlisting}
console.log(produk.nama);
console.log(produk.harga);
\end{lstlisting}

Menggunakan bracket notation:

\begin{lstlisting}
console.log(produk["kategori"]);
\end{lstlisting}

\subsection{Mengubah Properti Object}

Properti object dapat diubah atau ditambahkan.

\begin{lstlisting}
const produk = {
    nama: "Kopi Susu",
    harga: 15000
};

produk.harga = 18000;
produk.stok = 20;

console.log(produk);
\end{lstlisting}

\subsection{Array of Objects}

Dalam aplikasi nyata, data biasanya tidak hanya
terdiri dari satu object.

Kita dapat menggabungkan array dan object untuk
menyimpan banyak data produk.

\begin{lstlisting}
const daftarProduk = [
    {
        nama: "Kopi Susu",
        harga: 15000
    },
    {
        nama: "Teh Manis",
        harga: 8000
    },
    {
        nama: "Roti Bakar",
        harga: 12000
    }
];
\end{lstlisting}

Untuk mengakses nama produk pertama:

\begin{lstlisting}
console.log(daftarProduk[0].nama);
\end{lstlisting}

Untuk menampilkan seluruh produk:

\begin{lstlisting}
for (const produk of daftarProduk) {
    console.log(produk.nama);
}
\end{lstlisting}

Konsep array of objects akan sering digunakan
ketika kita mulai mengolah data produk pada
proyek website UMKM.

% ==================================================
\section{Praktik Terpadu}
% ==================================================

Pada bagian ini, kita akan menggabungkan beberapa
konsep yang telah dipelajari ke dalam satu program
sederhana.

\subsection{Studi Kasus: Sistem Perhitungan Belanja}

Buat sebuah file bernama \texttt{script.js}.

Program akan menyimpan beberapa data produk,
menghitung total harga, dan memberikan informasi
berdasarkan jumlah belanja.

\begin{lstlisting}
const daftarProduk = [
    {
        nama: "Kopi Susu",
        harga: 15000,
        jumlah: 2
    },
    {
        nama: "Roti Bakar",
        harga: 12000,
        jumlah: 1
    },
    {
        nama: "Teh Manis",
        harga: 8000,
        jumlah: 3
    }
];

function hitungTotal(harga, jumlah) {
    return harga * jumlah;
}

let totalBelanja = 0;

console.log("DAFTAR BELANJA");

for (const produk of daftarProduk) {
    const subtotal = hitungTotal(
        produk.harga,
        produk.jumlah
    );

    totalBelanja += subtotal;

    console.log(
        produk.nama + ": Rp" + subtotal
    );
}

console.log("----------------");

console.log(
    "Total: Rp" + totalBelanja
);

if (totalBelanja >= 50000) {
    console.log("Selamat, Anda mendapat diskon!");
} else {
    console.log("Belum memenuhi syarat diskon.");
}
\end{lstlisting}

\subsection{Analisis Program}

Program tersebut menggabungkan beberapa konsep:

\begin{enumerate}
    \item \textbf{Array of Objects}: menyimpan beberapa
    data produk.
    \item \textbf{Function}: menghitung subtotal
    berdasarkan harga dan jumlah.
    \item \textbf{Perulangan}: memproses setiap produk
    secara berurutan.
    \item \textbf{Variabel}: menyimpan total belanja.
    \item \textbf{Operator}: melakukan perhitungan.
    \item \textbf{Percabangan}: menentukan apakah
    pelanggan memperoleh diskon.
\end{enumerate}

Jalankan program melalui browser dengan menghubungkan
file JavaScript ke HTML, kemudian buka tab Console
untuk melihat hasilnya.

% ==================================================
\section{Latihan Mandiri}
% ==================================================

Kerjakan latihan berikut untuk menguji pemahaman
terhadap materi.

\subsection{Latihan 1: Variabel dan Operator}

Buat variabel berikut:

\begin{itemize}
    \item Nama toko.
    \item Nama produk.
    \item Harga produk.
    \item Jumlah produk.
\end{itemize}

Hitung total harga menggunakan operator aritmatika
dan tampilkan hasilnya menggunakan
\texttt{console.log()}.

\subsection{Latihan 2: Percabangan}

Buat variabel \texttt{stok} dengan nilai tertentu.

Gunakan percabangan untuk menentukan pesan:

\begin{itemize}
    \item Jika stok lebih dari nol, tampilkan
    \texttt{Produk tersedia}.
    \item Jika stok sama dengan nol, tampilkan
    \texttt{Produk habis}.
    \item Jika stok kurang dari nol, tampilkan
    \texttt{Data stok tidak valid}.
\end{itemize}

\subsection{Latihan 3: Perulangan}

Buat array yang berisi minimal lima nama produk.

Gunakan perulangan untuk menampilkan seluruh
nama produk ke console.

\subsection{Latihan 4: Fungsi}

Buat fungsi bernama \texttt{hitungDiskon()}.

Fungsi menerima dua parameter:

\begin{itemize}
    \item Harga awal.
    \item Persentase diskon.
\end{itemize}

Fungsi harus mengembalikan harga setelah diskon.

Contoh:

\begin{lstlisting}
hitungDiskon(100000, 10);
\end{lstlisting}

Hasil yang diharapkan:

\begin{lstlisting}
90000
\end{lstlisting}

\subsection{Latihan 5: Object dan Array}

Buat array bernama \texttt{daftarProduk} yang
berisi minimal tiga object produk.

Setiap object harus memiliki:

\begin{itemize}
    \item nama
    \item harga
    \item stok
\end{itemize}

Kemudian:

\begin{enumerate}
    \item Tampilkan seluruh nama produk.
    \item Tampilkan harga setiap produk.
    \item Hitung jumlah seluruh stok.
    \item Tampilkan produk yang stoknya lebih dari nol.
\end{enumerate}

% ==================================================
\section{Evaluasi Pemahaman}
% ==================================================

Jawablah pertanyaan berikut menggunakan pemahaman
sendiri. Peserta tidak cukup hanya menyalin kode
dari materi.

\begin{enumerate}
    \item Apa perbedaan penggunaan \texttt{let}
    dan \texttt{const}?

    \item Apa perbedaan tipe data String, Number,
    dan Boolean?

    \item Mengapa operator \texttt{===} lebih
    disarankan daripada \texttt{==}?

    \item Kapan kita menggunakan percabangan
    \texttt{if-else}?

    \item Apa perbedaan perulangan \texttt{for}
    dan \texttt{while}?

    \item Apa perbedaan parameter dan argument
    dalam sebuah fungsi?

    \item Mengapa array menggunakan index yang
    dimulai dari nol?

    \item Apa perbedaan array dan object?

    \item Mengapa array of objects berguna untuk
    menyimpan data produk?

    \item Jelaskan bagaimana program perhitungan
    belanja pada materi ini bekerja.
\end{enumerate}

% ==================================================
\section{Kriteria Penyelesaian}
% ==================================================

Peserta dinyatakan telah menyelesaikan materi ini
apabila mampu:

\begin{itemize}[label=$\square$]
    \item Membuat file JavaScript eksternal dan
    menghubungkannya dengan HTML.
    \item Menggunakan console untuk melihat hasil
    program dan membantu debugging.
    \item Membuat serta menggunakan variabel.
    \item Memahami tipe data dasar.
    \item Menggunakan operator aritmatika,
    perbandingan, dan logika.
    \item Membuat percabangan.
    \item Menggunakan perulangan.
    \item Membuat fungsi dengan parameter dan return.
    \item Mengolah array dan object sederhana.
    \item Menyelesaikan latihan tanpa hanya menyalin
    contoh dari materi.
\end{itemize}

% ==================================================
\section{Penutup}
% ==================================================

JavaScript memberikan kemampuan kepada sebuah
website untuk menjalankan logika dan merespons
kebutuhan pengguna.

Pada materi ini, kita telah mempelajari fondasi
pemrograman yang akan digunakan kembali dalam
materi-materi berikutnya.

Variabel membantu menyimpan data, operator
memproses nilai, percabangan mengambil keputusan,
perulangan mengerjakan tugas berulang, sedangkan
fungsi membantu menyusun kode agar lebih terstruktur.

Array dan object kemudian memberikan cara untuk
mengelola kumpulan data yang lebih kompleks.

Konsep-konsep tersebut merupakan fondasi penting
sebelum kita mempelajari bagaimana JavaScript
berinteraksi langsung dengan elemen HTML melalui
DOM.

Ingat, memahami logika lebih penting daripada
sekadar menghafal sintaks.

\begin{center}
    \vspace{0.5cm}
    \textbf{\Large Learn. Build. Understand.}

    \vspace{0.2cm}
    \textit{PROJECT DAWET}
\end{center}

% ==================================================
\section*{Referensi Pembelajaran}
% ==================================================

Gunakan dokumentasi berikut sebagai referensi
tambahan:

\begin{enumerate}
    \item MDN Web Docs -- JavaScript:
    \url{https://developer.mozilla.org/en-US/docs/Web/JavaScript}

    \item MDN -- JavaScript First Steps:
    \url{https://developer.mozilla.org/en-US/docs/Learn_web_development/Core/Scripting}

    \item MDN -- JavaScript Guide:
    \url{https://developer.mozilla.org/en-US/docs/Web/JavaScript/Guide}

    \item JavaScript.info:
    \url{https://javascript.info/}
\end{enumerate}

\end{document}
```